\section{Conclusion}
\label{sec:conclusion}

On a controlled two-dimensional model with a quadrature-exact reference free
energy, a Fisher--Rao birth--death correction can help the Adaptive Biasing
Force method, but only as a \emph{carefully controlled marginal reallocation
mechanism}. The practical, estimated-target FR improves the time-integrated
free-energy error by $\EstIntegImpPct\%$ over ABF only at its best configuration
and gives smaller final-budget gains; it is best understood here as a
convergence accelerator and stabiliser rather than a qualitative rescue. The
improvement is not unconditional --- it depends on the FR rate, bandwidth and
burn-in, and poor choices match or underperform ABF --- and across all
successful settings FR is \emph{weak} (event fractions of order $10^{-5}$),
\emph{target-consistent}, and does not visibly corrupt the conditional
distribution $Y\mid X=x$ on which ABF depends.

Two interpretive points are essential. First, the practical value of FR is
governed by the accuracy and stability of the online target free-energy
estimate: the gap between the estimated target ($\EstIntegImpPct\%$
integrated-error reduction) and the oracle control ($\OracIntegImpPct\%$)
localises the bottleneck in target estimation, not in the birth--death
mechanism. Second, the oracle-target experiments must be read \emph{only} as
diagnostic evidence that the mechanism has potential when the target marginal
direction is accurate; because the oracle uses the reference free energy
$\Fref$, it is not a deployable algorithm and is never recommended as one.
Whether FR is ``icing on the cake'' or ``snow in the winter'' will be settled
not on this easy example but in the harder regimes outlined in
\cref{sec:discussion}.
